\begin{alist}
\item Οι πιθανές ακέραιες ρίζες της εξίσωσης βρίσκονται στους διαιρέτες του $-1$ και είναι οι αριθμοί $\pm 1$. \\
Αλλά $f(1) = 3 \ne 0$ και $f(-1) = -3 \ne 0$. \\
Επομένως η εξίσωση $f(x) = 0$ δεν έχει ακέραιες ρίζες.

\item 
\begin{rlist}
\item Η εξίσωση $f(x) = 0$ έχει μία ρίζα, διότι η γραφική παράσταση της συνάρτησης $f$ τέμνει τον άξονα $x'x$ στο σημείο $\text{A}$.
\item Είναι $f(1) = 3 > 0$ και $f(0) = -1 < 0$. \\
Επομένως, σύμφωνα με το θεώρημα που προσδιορίζει τη ρίζα μιας πολυωνυμικής συνάρτησης με προσέγγιση, υπάρχει μία τουλάχιστον ρίζα της εξίσωσης $f(x) = 0$ μεταξύ των αριθμών $0$ και $1$.
\end{rlist}
\end{alist}
